\chapter{Implementasi Orkestrasi Apache Airflow}

% ============================================================
%  BAB V: IMPLEMENTASI ORKESTRASI APACHE AIRFLOW
% ============================================================
 
\section{Arsitektur DAG (Directed Acyclic Graph)}
 
Apache Airflow bertindak sebagai mesin pengatur dan penjadwal utama (\textit{enterprise orchestrator}) untuk mengotomatisasi seluruh siklus data dalam sistem DWBI rumah sakit ini. Di dalam Airflow, rangkaian alur kerja didefinisikan sebagai \textbf{Directed Acyclic Graph (DAG)} bernama \texttt{02\_healthcare\_data\_pipeline}. Karakteristik \textit{"directed"} menandakan alur eksekusi memiliki arah ketergantungan yang jelas dari awal hingga akhir, sedangkan \textit{"acyclic"} menjamin bahwa tidak ada tugas yang berjalan berulang atau membentuk putaran tak berujung (\textit{infinite loop}).
 
Konfigurasi parameter operasional dan manajemen risiko dari DAG ini dirancang menggunakan argumen bawaan (\textit{default args}) sebagai berikut:
 
\begin{itemize}
    \item \textbf{Owner (\texttt{kelompok\_healthcare}):} Digunakan sebagai identitas pengelola aset \textit{pipeline} untuk mempermudah kategorisasi, audit log, dan pemfilteran pada dasbor Web UI Airflow.
    \item \textbf{Depends on Past (False):} Keberhasilan eksekusi \textit{pipeline} hari ini tidak digantungkan pada status keberhasilan eksekusi hari sebelumnya, sehingga mencegah kemacetan massal (\textit{data blocking}).
    \item \textbf{Retries (1) \& Retry Delay (2 menit):} Jika sebuah \textit{task} gagal akibat masalah jaringan atau database PostgreSQL yang sedang sibuk, Airflow akan menjadwalkan ulang eksekusi otomatis satu kali lagi setelah jeda 2 menit.
    \item \textbf{Schedule Interval (None) \& Catchup (False):} \textit{Pipeline} dikonfigurasi untuk berjalan berdasarkan pemicu manual (\textit{Manual Trigger}). Penonaktifan \textit{catchup} mencegah Airflow melakukan eksekusi susulan (\textit{backfilling}) terhadap jadwal historis masa lalu.
\end{itemize}
 
\section{Topologi dan Penjelasan Task Dependencies}
 
\textit{Pipeline} dipecah menjadi 4 \textit{task} diskrit yang dikelola oleh dua operator utama: \texttt{BashOperator} (untuk mengeksekusi skrip eksternal di lingkungan OS) dan \texttt{PythonOperator} (untuk mengeksekusi fungsi di dalam memori Python). Keempat \textit{task} diikat menggunakan operator dependensi sekuensial linear (\texttt{>>}), membentuk urutan kronologis sebagai berikut:
 
\textbf{Task 1: \texttt{run\_data\_generator} (BashOperator)}
 
Airflow memicu lingkungan \textit{shell} Linux untuk menjalankan skrip pembuat data sintetik.
 
\begin{lstlisting}[language=Python, caption=Definisi Task run\_data\_generator]
task_generate_data = BashOperator(
    task_id='run_data_generator',
    bash_command='python /opt/airflow/data_generator/src/generator.py',
)
\end{lstlisting}
 
Airflow memastikan file \texttt{generator.py} selesai dieksekusi 100\% dan menghasilkan 3 file CSV di folder \textit{shared storage} sebelum mengizinkan aliran data berpindah ke tahap berikutnya.
 
\textbf{Task 2: \texttt{ingest\_csv\_to\_database} (PythonOperator)}
 
Airflow memanggil fungsi Python internal untuk memindahkan data dari file \texttt{.csv} ke tabel relasional database.
 
\begin{lstlisting}[language=Python, caption=Definisi Task ingest\_csv\_to\_database]
task_ingest_data = PythonOperator(
    task_id='ingest_csv_to_database',
    python_callable=ingest_csv_to_postgres,
)
\end{lstlisting}
 
Fungsi \texttt{ingest\_csv\_to\_postgres} menggunakan \textit{connection string} PostgreSQL \texttt{postgresql+psycopg2:\\//airflow:airflow@postgres:5432/airflow} untuk membaca CSV via \textit{Dataframe} Pandas dan menuliskannya ke tabel \texttt{src\_pasien}, \texttt{src\_dokter}, dan \texttt{src\_kunjungan} dengan metode \textit{overwrite} (\texttt{if\_exists='replace'}).
 
\textbf{Task 3: \texttt{dbt\_run\_transformation} (BashOperator)}
 
Airflow memerintahkan \texttt{dbt} sebagai \textit{third-party tool} untuk memproses transformasi.
 
\begin{lstlisting}[language=Python, caption=Definisi Task dbt\_run\_transformation]
task_dbt_run = BashOperator(
    task_id='dbt_run_transformation',
    bash_command=(
        'cd /opt/airflow/dbt_project && '
        'dbt run --target docker_env --profiles-dir .'
    ),
)
\end{lstlisting}
 
Airflow mengirimkan sinyal perintah CLI \texttt{dbt run} untuk memproses pembentukan \textit{Star Schema} (Tabel Fakta dan Dimensi) di dalam database.
 
\textbf{Task 4: \texttt{dbt\_data\_quality\_test} (BashOperator)}
 
Airflow mengeksekusi proses pengujian kualitas data sebelum laporan siap dikonsumsi.
 
\begin{lstlisting}[language=Python, caption=Definisi Task dbt\_data\_quality\_test]
task_dbt_test = BashOperator(
    task_id='dbt_data_quality_test',
    bash_command=(
        'cd /opt/airflow/dbt_project && '
        'dbt test --target docker_env --profiles-dir .'
    ),
)
\end{lstlisting}
 
Jika pengujian integritas (seperti deteksi \textit{null value} atau duplikasi kunci utama) menemukan kegagalan, Airflow akan menandai \textit{task} ini berwarna merah (\textit{Failed}) dan menghentikan seluruh \textit{pipeline}.
 
\section{Panduan Dokumentasi Antarmuka Airflow UI}
 
Untuk membuktikan bahwa arsitektur orkestrasi ini berjalan dengan sukses di lingkungan Docker, lampirkan beberapa tangkapan layar dari Apache Airflow Web UI (akses via \texttt{localhost:8080}) dengan penjelasan sebagai berikut:
 
\begin{itemize}
    \item \textbf{Screenshot Graph View:} 
   \begin{center}
        \centering
        \includegraphics[width=0.9\textwidth]{"./figures/5-1_DAG_graph.png"}
        \captionof{figure}{Gambar 5.1 Tampilan Graph View Ketergantungan Task.}
        \label{fig:5.1.DAG_graph}
    \end{center}
    
    Gambar ini mendokumentasikan visualisasi topologi linear dari pipeline. Terlihat jelas arah panah eksekusi dari kiri ke kanan yang dimulai dari pembuatan data hingga pengujian kualitas data. Seluruh bingkai kotak (node) harus dipastikan berwarna hijau tua (dark green), yang menjadi indikator absolut di Airflow bahwa seluruh perintah script sukses dieksekusi tanpa interupsi.
 
    \item \textbf{Screenshot Grid / Tree View:} 
    \begin{center}
        \centering
        \includegraphics[width=0.4\textwidth]{"./figures/5-2_DAG_grid_time.png"}
        \captionof{figure}{Gambar 5.2 Tampilan Grid View Histori Eksekusi Pipeline.}
        \label{fig:5.2.DAG_grid_time}
    \end{center}
    
    Dokumentasi ini memperlihatkan rekam jejak (log run) dari eksekusi DAG. Lajur kotak vertikal berwarna hijau penuh menandakan keberhasilan run secara utuh pada baris waktu tertentu. Melalui tampilan ini, pengelola Data Warehouse dapat memantau durasi performa kecepatan eksekusi tiap task dan mendeteksi jika terjadi lonjakan beban komputasi (bottleneck).
 
    \item \textbf{Screenshot Task Log:}
    \begin{center}
        \centering
        \includegraphics[width=0.9\textwidth]{"./figures/5-3_DAG_task_log.png"}
        \captionof{figure}{Gambar 5.3 Tampilan Log Output Terminal Eksekusi Task.}
        \label{fig:5.3.DAG_task_log}
    \end{center}
    
    Dokumen log menunjukkan status keluaran sistem yang bersih dengan kode balik \texttt{INFO - Command exited with return code 0} pada task \texttt{run\_data\_generator}, membuktikan Airflow berhasil menyelesaikan pemanggilan fungsi eksternal dengan sukses.
\end{itemize}

